% Part: first-order-logic
% Chapter: syntax-and-semantics
% Section: introduction

\documentclass[../../../include/open-logic-section]{subfiles}

\begin{document}

\olfileid[hi]{fol}{syn}{int}

\olsection{परिचय}

प्रथम-क्रम तर्क के सिद्धांत और अधिसिद्धांत को विकसित करने के लिए हमें पहले
उसकी अभिव्यक्तियों के वाक्यविन्यास और अर्थविज्ञान को परिभाषित करना होगा।
प्रथम-क्रम तर्क की अभिव्यक्तियाँ पद और !!{formula}s हैं। पद !!{variable}s,
!!{constant}s और !!{function}s से बनते हैं। इसके बाद !!^{formula}s,
!!{predicate}s और पदों से बनते हैं (इनसे सबसे छोटे, ``परमाण्विक''
!!{formula}s बनते हैं); फिर तार्किक संयोजकों और परिमाणकों का उपयोग करके
परमाण्विक !!{formula}s से अधिक जटिल सूत्र बनाए जा सकते हैं। निर्माण-नियम
देने के अनेक ढंग हैं; यहाँ हम केवल एक संभव ढंग देते हैं। अन्य पद्धतियाँ भिन्न
प्रतीक चुनेंगी, संयोजकों के भिन्न समुच्चयों को आदिम मानेंगी और कोष्ठकों का
भिन्न प्रकार से उपयोग करेंगी (या तथाकथित पोलिश संकेतन की तरह उनका बिल्कुल
उपयोग नहीं करेंगी)। फिर भी सभी दृष्टिकोणों में यह बात समान है कि निर्माण-नियम
पदों और !!{formula}s के समुच्चय को \emph{आगमनात्मक रूप से} परिभाषित करते हैं।
परिभाषा ठीक ढंग से दी गई हो तो निर्माण-नियमों के अनुसार प्रत्येक अभिव्यक्ति
मूलतः केवल एक ही प्रकार से बन सकती है। अभिव्यक्तियों को \emph{अद्वितीय रूप से
पठनीय} बनाने वाली आगमनात्मक परिभाषा के कारण हम उसी विधि---आगमनात्मक
परिभाषा---से इन अभिव्यक्तियों को अर्थ दे सकते हैं।

अभिव्यक्तियों को अर्थ देना अर्थविज्ञान का विषय है। अर्थविज्ञान की केंद्रीय
धारणा किसी !!a{structure} में संतुष्टि की है। !!^a{structure} भाषा के
निर्माण-अवयवों को अर्थ देती है: !!a{domain} वस्तुओं का एक अरिक्त समुच्चय
है। परिमाणकों का निर्वचन इस प्रांत पर विचरण करने के रूप में किया जाता है,
!!{constant}s को प्रांत के अवयव दिए जाते हैं, !!{function}s को !!{domain}
से उसी में जाने वाले फलन दिए जाते हैं, और !!{predicate}s को !!{domain} पर
संबंध दिए जाते हैं। मूल शब्दावली के इन नियतनों के साथ !!{domain} मिलकर
!!a{structure} बनाता है। !!^{variable}s, !!{formula}s में आ सकते हैं; अतः
अर्थविज्ञान देने के लिए हमें उन्हें भी !!{domain} के !!{element}s नियत करने
होते हैं---इसे चर-नियतन कहते हैं। अंततः संतुष्टि-संबंध इन सबको एक साथ जोड़ता
है। कोई !!^a{formula}, !!a{structure}~$\Struct{M}$ में चर-नियतन~$s$ के
सापेक्ष संतुष्ट हो सकता है; इसे $\Sat{M}{!A}[s]$ लिखते हैं। इस संबंध की
परिभाषा भी~$!A$ की संरचना पर आगमन से दी जाती है। उदाहरणतः तार्किक संयोजकों
की सत्यता-सारणियों का उपयोग करके $!A \land !B$ की संतुष्टि को $!A$ और~$!B$
की संतुष्टि (या असंतुष्टि) के आधार पर परिभाषित किया जाता है। फिर पता चलता है
कि यदि !!{formula}~$!A$ कोई !!a{sentence} है, अर्थात् उसमें कोई मुक्त चर
नहीं है, तो चर-नियतन अप्रासंगिक है। इसलिए हम सीधे कह सकते हैं कि कोई
!!{sentence}, किसी !!{structure} में संतुष्ट है (या नहीं)।

!!{sentence}s के लिए संतुष्टि-संबंध $\Sat{M}{!A}$ के आधार पर हम वैधता,
अर्थगत तात्पर्य और संतुष्टनीयता की मूल अर्थगत धारणाएँ परिभाषित कर सकते हैं।
कोई वाक्य वैध है, $\Entails !A$, यदि प्रत्येक संरचना उसे संतुष्ट करती है।
!!{sentence}s का कोई समुच्चय उससे तात्पर्य रखता है, $\Gamma \Entails !A$,
यदि~$\Gamma$ के सभी !!{sentence}s को संतुष्ट करने वाली प्रत्येक
!!{structure},~$!A$ को भी संतुष्ट करती है। और वाक्यों का कोई समुच्चय
संतुष्टनीय है यदि कोई !!{structure} उसके सभी !!{sentence}s को एक साथ संतुष्ट
करती है। चूँकि !!{formula}s आगमनात्मक रूप से परिभाषित हैं और संतुष्टि भी
!!{formula}s की संरचना पर आगमन से परिभाषित होती है, इसलिए हम अपने अर्थविज्ञान
के गुण सिद्ध करने तथा परिभाषित अर्थगत धारणाओं को परस्पर संबंधित करने के लिए
आगमन का उपयोग कर सकते हैं।




\end{document}
